\chapter{Introduction}\label{ch:introduction}

\draftnote{Motivate the work, state the research questions, list contributions, and outline the rest of the thesis.}

\section{Motivation}\label{sec:motivation}

\draftnote{Chess engines on embedded devices (Wio Terminal: 192\,KB RAM, 500\,KB flash) need an efficient evaluation function such as NNUE. Handcrafted bucketing (piece count, queen presence) is heuristic and may not capture the true structure of the state space.}

\section{Problem Statement}\label{sec:problem-statement}

\draftnote{How can we automatically learn a partition of the state space that maximises expert specialisation, and how can we do this without human supervision?}

\section{Contributions}\label{sec:contributions}

\draftnote{(i)~An unsupervised bucketing method based on sample gradients (gradients with respect to head parameters). (ii)~A lightweight dispatcher that routes positions to the right expert at inference time. (iii)~A Mixture of Experts NNUE architecture for resource-constrained devices. (iv)~Empirical validation on chess against a single-head baseline.}

\section{Thesis Outline}\label{sec:outline}

\draftnote{Roadmap: Chapter~\ref{ch:background} reviews engines, NNUE, MoE, sample gradients, and bucketing. Chapter~\ref{ch:method} presents the five-step unsupervised bucketing pipeline. Chapter~\ref{ch:implementation} describes the chess/NNUE instantiation. Chapter~\ref{ch:experiments} reports experiments. Chapter~\ref{ch:discussion} discusses results and limitations. Chapter~\ref{ch:conclusion} concludes and sketches future work.}
